\paragraph{}Systemy rozpoznawania tożsamości na podstawie obrazu twarzy osiągają imponujące wyniki w warunkach kontrolowanych dzięki nowoczesnym głębokim sieciom neuronowym i funkcjom strat opartym na marginesach kątowych, takich jak ArcFace. Ich odporność jednak zazwyczaj spada w obecności okluzji - kluczowego wyzwania w rzeczywistych aplikacjach, w tym systemach monitoringu, biometrii bezpieczeństwa i urządzeniach mobilnych. Niniejsza praca analizuje wpływ częściowej okluzji twarzy, szczególnie zasłonięcia obszaru oczu, na dokładność identyfikacji osób oraz ocenia skuteczność modyfikacji architektur zwiększających odporność na tego typu zakłócenia.

Implementujemy i ewaluujemy trzy warianty architektur opartych na IResNet50: model bazowy, model z rozproszoną uwagą CBAM w każdym bloku rezydualnym oraz model z pomocniczą gałęzią Transformer. Ewaluacja na testowym zbiorze CASIA-WebFace z deterministyczną okluzją 20px wskazuje, że rozproszona uwaga poprawia odporność na zasłonięcie: CBAM\_Block osiąga 90.16\% Rank-1, przewyższając baseline (87.11\%) o 3.05 pp. Wariant Transformer uzyskuje 85.00\% Rank-1.

Analiza potwierdza, że hierarchiczna filtracja zakłóceń na wielu poziomach abstrakcji cech jest skutecznym kierunkiem zwiększania odporności modeli rozpoznawania twarzy w scenariuszach z okluzją. Proponowane modyfikacje oparte na CBAM zachowują prostotę wdrożenia przy jednoczesnej poprawie jakości identyfikacji.
